%% kapa: exam=1 round=30 year=2019 period=1 form=A subject=real_estate subject_ko=부동산학원론 seq=40 number=120
%% source: Q-Net 감정평가사 기출문제 (cst00309) artlSeq=5205205 file=extracted/1008/5205205/2019년 제30회 감정평가사 1차시험 문제지/2019년 제30회 감정평가사 1차시험 문제지/2019년도 제30회 감정평가사 1차 1교시 A형.hwp
%% conversion: kordoc 4.24.51 -> tools/gen1.py (auto); review=proofread-sonnet; flags=none
\begin{question}{120}
자산 A, B, C에 대한 경제상황별 예상수익률이 다음과 같을 때, 이에 관한 설명으로 옳지 않은 것은? (단, 호황과 불황의 확률은 같음) \par
{\small\centering
\begin{tabular}{|c|c|c|}
\hline
\multirow{2}{*}{\makecell{구 분}} & \multicolumn{2}{|c|}{\makecell{경제상황별 예상수익률(\%)}} \\
\cline{2-3}
 & \makecell{호황} & \makecell{불황} \\
\hline
\makecell{자산A} & \makecell{8} & \makecell{4} \\
\hline
\makecell{자산B} & \makecell{12} & \makecell{8} \\
\hline
\makecell{자산C} & \makecell{16} & \makecell{10} \\
\hline
\end{tabular}\par}

\begin{choices}
  \item 기대수익률은 자산C가 가장 높고, 자산A가 가장 낮다.
  \item 합리적인 투자자라면 자산A와 자산B 중에서는 자산B를 투자안으로 선택한다.
  \item 평균분산지배원리에 따르면 자산C가 자산B를 지배한다.
  \item 자산B의 변동계수(Coefficient of variation)는 0.2이다.
  \item 자산C가 상대적으로 다른 자산에 비해서 위험이 높다.
\end{choices}
\end{question}
\answer{3}
